\section{Algorytm}

Filtr Savitzky-Golay pozwala wygładzić cyfrowy sygnał. Jego działanie opiera się na przybliżaniu lokalnych zestawów sąsiadujących ze sobą punktów wielomianami niskiego rzędu.

Mając dany sygnał x[n], poszukujemy wielomianu
\begin{equation}
p(n) = \sum\limits_{k=0}^N a_k n^k
\end{equation}
który w otoczeniu $2M+1$ punktów minimalizuje kwadrat błędu aproksymacji
\begin{equation}
\epsilon_N = \sum\limits_{n=-M}^M (p(n) - x[n])^2
\end{equation}



W praktyce filtr Savitzky-Golay realizuje swoje zadanie obliczając splot kilku aktualnie branych pod uwagę punktów z aproksymacją odpowiedzi impulsowej ...


Przebieg filtracji uzyskanej za pomocą prototypu Python przedstawia wykres \ref{rys:savitzky_py}


Automatyczna numeracja obrazków: obrazek \ref{rys:wzor}.

\begin{figure}[!htb]
  \begin{center}
    \includegraphics[width=12cm,trim=1.6cm 6.9cm 1.7cm 8.8cm,clip]
    {img/test.jpg}
  \end{center}
  \caption{Obrazek testowy}
  \label{rys:savitzky_py}
\end{figure}


